\centerline{\Huge \bf  Тренировочная Олимпиада 3}
\centerline{\Huge \bf 7 класс}
\centerline{\Huge \bf Уровень: Регион Эйлера. День 2.}

\begin{flushright}
    \scriptsize{Глава 3. День, когда Токио-3 замер.}
\end{flushright}


\problem{1}
На день рождения Расула Владимировича пришло 49 учеников ЭлМШ. Известно, что любые 32 ребёнка суммарно съели хотя бы половину от всей еды. Какую наибольшую часть от еды смог бы съесть один Эрдем? Расул Владимирович к еде не прикасался, а дети по итогу съели всю имеющуюся еду.
\begin{flushright}
    \textit{\small{(Гасанов Р.В.)}}
\end{flushright}

\problem{2}
На полке у Максима Алексеевича лежит в некотором порядке 11 томов манги Человек-бензопила. Докажите, что их можно расставить в правильном порядке, если разрешается выполнять следующие действия:

$-$ Поменять местами первые две манги

$-$ Положить самую правую мангу в начало
\begin{flushright}
    \textit{\small{(Гасанов Р.В.)}}
\end{flushright}

\problem{3}
Катет $MK$ прямоугольного треугольника $MNK$ продлили за точку $K$ на отрезок $KT$, так что треугольники $MNK$ и $TNK$ равновелики. Точка $E$ середина отрезка $NT$ и $ME$ пересекается с $NK$ в точке $L$. Найдите отношение площади треугольника $KET$ к площади треугольника $MNL$.
\begin{flushright}
    \textit{\small{(Доржеев А.А.)}}
\end{flushright}

\problem{4}     
Сколькими способами, меняя между собой некоторые буквы (не обязательно различные), из слова БУСЯДУСЯЛЮСЯМАРУСЯКРИСТИВАНГА можно получить другое слово, которое тоже означало бы все имена кошек Расула Владимировича.
\begin{flushright}
    \textit{\small{(Гасанов Р.В.)}}
\end{flushright}

\problem{5} 
Муравьиная колония насчитывает 20 гнёзд. Муравьиные Королевы отдали приказ, соединить гнезда так, что бы из любого гнезда можно попасть в любой другой. Но проблема вот в чём, земля достаточно рыхлая, и невозможно построить более 4-х туннелей в одном гнезде. Муравьи истинные трудяги, и по этому стараются выполнять любую работу как можно быстрее, соответственно проходить кратчайшим путями. Помогите Билданту понять, какой протяженности самый длинный путь из одного гнезда в другой, если передвигаться по муравьиному ( по наименьшему количеству туннелей)?
\begin{flushright}
    \textit{\small{(Очиров O.В.)}}
\end{flushright}
